\section{Related Work}
\label{sec:related}

\subsection{Statecharts and State-Machine Notations}

Harel's statecharts~\cite{harel1987statecharts} introduced hierarchical,
concurrent finite state machines as a visual formalism for reactive systems,
providing the first notation expressive enough to describe realistic
event-driven behavior without combinatorial explosion. The influence on Trenza
is direct: contexts and transitions are semantically grounded in the same
intuition that state topology is the right primitive for behavioral
specification. The departure is equally direct. Statecharts are a diagrammatic
notation; their primary artifact is a drawing, not a text file---incompatible
with version control, diff-based review, and consumption by language models.
Trenza inverts the artifact hierarchy: the \texttt{.trz} specification is the
source of truth stored as plain text under Git, and the Mermaid statechart is a
derived projection synthesized at compile time.

SCXML~\cite{scxml2015} and XState~\cite{xstate2019} bring statecharts to
software environments as XML vocabularies and JavaScript libraries
respectively. Both remain implementation-level formalisms: they specify how a
runtime should behave, not what behavioral contracts hold across all contexts.
Neither enforces role exhaustiveness, data-scope least privilege, or
multi-target synthesis.

\subsection{Conceptual Antecedents}

Reenskaug and Coplien's Data-Context-Interaction architecture~\cite{reenskaug2009dci}
elevated roles to first-class citizens by separating the stable structure of
domain objects (Data) from situational use-case logic (Context) and dynamic
execution paths (Interaction). Trenza inherits DCI's core insight directly: a
datum is inert until it assumes a transient role within a named context; outside
that context, the role and its permissions do not exist. Where Trenza departs
from DCI is in formalizing this insight as a set of mechanically enforced rules.
DCI is an architectural philosophy; it imposes no compiler-checkable constraints
on which roles must appear in which contexts. Trenza does: Rule~1 (Completeness)
requires every role-event pair to be handled in every context; Rule~6 (Data
conformance) enforces access-scope least privilege structurally.

Pawson's Naked Objects~\cite{pawson2002nakedobj} argues that domain behavior
should be exposed directly through the user interface, without intermediate
presentation layers. The underlying commitment is shared: behavior should be
visible at the level where decisions are taken. In Naked Objects, the UI exposes
what the domain object does; in Trenza, the \texttt{.trz} exposes what every
(role, event) pair does in every context. Both proposals attack the same
defect---behavior hidden behind a layer of indirection---at different points in
the stack.

NetKernel~\cite{rodgers2006netkernel} introduced Resource-Oriented Computing,
in which computation is described as a sequence of named, addressable,
side-effect-bounded transformations. Trenza shares the intuition that naming is
the right lever for making systems tractable---specifically, naming contexts as
the addressable unit of behavioral specification. The lineage is indirect and is
declared as such.

\subsection{Formal Specification Languages}

TLA+~\cite{lamport2002tla} and Alloy~\cite{jackson2006alloy} are the most
capable lightweight formal methods for software specification. TLA+ can express
and model-check liveness and safety properties of distributed systems; Alloy
can find counterexamples to relational invariants across bounded domains. The
power of both tools is not in question. The constraint is access. Writing a
TLA+ specification requires comfort with temporal logic; writing an Alloy model
requires fluency with relational algebra. Neither was designed with LLM
consumption as a design goal. Trenza occupies a different point in the design
space: it sacrifices the expressive power of full temporal logic in exchange for
a constraint set that is simultaneously enforceable by a fast compiler and
traversable by an LLM as a finite truth table. The verification question in
Trenza is not ``does this system satisfy a temporal property under all
interleavings?'' but ``does this implementation correspond, entry by entry, to
what the specification declares?''---a weaker question that a language model
answers deductively rather than probabilistically.

\subsection{Model-Driven Engineering}

Executable UML~\cite{mellor2002xuml} and the broader MDA tradition pursued the
same fundamental goal as Trenza: elevate specification above implementation and
generate code from formal models. This project failed in practice for
compounding reasons~\cite{mernik2005when}: modeling notations grew to
accommodate every possible use case; artifacts lived in proprietary binary
formats incompatible with version control; code generation was partial, breaking
the round-trip guarantee and leaving models perpetually out of sync with
production code. Most consequentially for the present work, these models were
designed to be read by humans pointing at a GUI, not parsed or critiqued by a
language model. Trenza is designed for the opposite constraints: a single
plain-text file per specification, stored in Git, structured so that every
semantic contract appears exactly once, and processable by a compiler and an LLM
with equal facility.

\subsection{Language Design for Non-Expert Readers}

Hermans documents~\cite{hermans2021brain} that the design of programming
languages is a problem of working memory and chunking: languages that minimize
the number of unrelated concepts a reader must hold in mind simultaneously are
easier to learn and easier to review. Trenza's eight rules are deliberately
phrased so that each rule corresponds to a single chunk a reader can hold
(\emph{``every role handles every event in every context''}); the verifier's
diagnostics report violations in the same vocabulary. The design pressure that
produced this---the LLM must be able to read the rule and reason about it
without unpacking a definition---is the same pressure Hermans documents for
human readers.

Hedy~\cite{hermans2020hedy} demonstrates that a programming language's syntactic
surface can be designed to grow with the learner, with each level complete in
itself and each transition small enough to avoid overwhelming the reader.
Trenza's design takes a related stance for an LLM audience: the surface is
intentionally small and self-describing, so that a model with a \texttt{.trz}
file in context can reason about modifications without consulting external
documentation.

\subsection{DSLs for Reactive and State-Managed Systems}

The Elm architecture~\cite{czaplicki2012elm} and its descendants---including
typed reactive frameworks for other functional ecosystems---demonstrate that a
sufficiently restrictive type system, combined with a canonical state management
pattern, can eliminate entire classes of runtime errors. The design philosophy
is shared with Trenza: the right constraints imposed early reduce defect surface
area more effectively than post-hoc testing applied to an unconstrained language.
The difference lies in scope and synthesis target. Elm and similar systems are
single-paradigm, single-target: specification and implementation are the same
artifact. Trenza separates them explicitly. The same \texttt{.trz} specification
currently produces four coherent artifacts---Rust, TypeScript, Mermaid
schematics, and audit reports---from a single source. The completeness,
determinism, and reachability guarantees are enforced once, at specification
level, and hold across all targets simultaneously.

\subsection{Contemporary Systems}

Multi-agent coordination frameworks such as \texttt{claude-flow}
\cite{claudeflow2024} address a legitimate problem: distributing independent
workstreams across
specialized agents when no single model session can hold the full context.
That problem is real. The state dispersal problem described in
Section~\ref{sec:motivation} is different in kind. It is not a question of
cognitive capacity; it is the absence of a shared, formal account of what the
system's state is. Adding agents without adding a state model produces a system
that is more complex without being more coherent. The two approaches address
different failure modes and are not in competition for the same problem space.

The Model Context Protocol~\cite{anthropic2024mcp} standardizes the protocol
layer by which an LLM session communicates with external tools and resources.
Trenza occupies a different layer: it specifies \emph{what the system does},
not how the model invokes external capabilities. The two compose naturally---a
\texttt{.trz} file can declare \texttt{external:} modules whose runtime binding
is provided by MCP servers---but Trenza adds no constraints on the protocol
itself.
